\documentclass[11pt]{article}
\usepackage[margin=0.75in]{geometry}
\usepackage{booktabs}
\usepackage{amsmath}
\usepackage[hidelinks]{hyperref}
\setlength{\parindent}{0pt}
\setlength{\parskip}{4pt}
\setlength{\tabcolsep}{7pt}

\begin{document}
\begin{center}
{\Large Wildfire sensitivity at arrival: working note}\\[4pt]
Until-stable run 3110691 $\cdot$ 19 September 2026 $\cdot$ two rounds, about 63 minutes
\end{center}

\textbf{Current result.}
The original-order screening labels \textbf{T and CO sensitive}, and leaves
\textbf{LTS and PREC undetermined}. Eight tested variables enter the insensitive set:
PBLH, CFLOW, CFMID, SWGDN, LWP, OMEGA, LWC, and RH. Together with the ten starting
assumptions (P, LAND, SLP, WS, TS, DUST, EPV, CHL, DMS, SO2EM), that gives
18 insensitive variables. The starting ten were not tested. These are screening
labels, not established causal effects. We use only until-stable experiments for
reported results.

\textbf{What is fitted and held out?}
For each candidate $c$, a global-latent LSTM engressor learns
$c_{\mathrm{arrival}}=f(M_{1:241},\varepsilon)$ on clean trajectories, where $M$ is
the current insensitive input set. Inputs have shape $(N,241,|M|)$; training and
validation targets have shape $(N,1)$, using the final array position as arrival.
The model predicts one candidate value per draw. Recalibration and coverage also
use only that arrival value. We pool the original clean
training and test rows, then form four separate sets:

\begin{center}
\begin{tabular}{lrl}
\toprule
Clean set & Trajectories & Purpose \\
\midrule
Training & 23,564 & Fit the network \\
Validation & 3,110 & Early stopping; select best checkpoint \\
Recalibration & 1,350 & Fit the quantile correction \\
Control / test & 2,159 & Check corrected coverage \\
\bottomrule
\end{tabular}
\end{center}

The clean pool has 54,026 trajectories; the other 23,843 are unused after time-gap
purging and month matching. Recalibration and control are separately month-matched
to 6,786 wildfire trajectories. Splits use 30-day arrival-time blocks with history
gaps at both boundaries. Wildfire outcomes are used only for evaluation.

\textbf{What recalibration changes.}
Kuleshov, Fenner, and Ermon~\cite{kuleshov2018} fit a monotone map from predicted
CDF probabilities to observed frequencies, recommending isotonic regression.
Our implementation uses a simpler empirical quantile map. On the recalibration
rows, compute each observation's position $u_i=F_i(y_i)$ in the base model's
predictive distribution, approximated by randomized ranks among 400 draws.
Pool these positions across arrival observations, then use
\[
 p_q=\operatorname{Quantile}_q(\{u_i\}),
 \qquad Q_{\mathrm{corrected},x}(q)=Q_{\mathrm{base},x}(p_q).
\]
The network is unchanged. One mapping per candidate adjusts all its quantile
predictions. A central 90\% interval uses corrected 5th and 95th percentiles;
a 95\% interval uses corrected 2.5th and 97.5th percentiles. The same mapping is
applied to control and wildfire data. Fitting the correction does not guarantee
that it achieves the target coverage on the separate control set.

\textbf{Decision rule.}
A model is accepted if its control coverage at both 90\% and 95\% lies within
$\pm3$ estimated standard errors of nominal. Standard errors use 30-day blocks
and the same observation weighting as pooled coverage. A failed check leaves
$c$ undetermined and outside $M$. For accepted models, a control-minus-wildfire
95\% coverage drop greater than two percentage points is sensitive; otherwise
$c$ enters $M$. Sensitive and undetermined candidates are retried after $M$ grows.
The run stops when a round adds nobody.

\newpage
\small
\textbf{Latest assessment of each tested candidate.}
Original-order run 3110691: arrival coverage in percent; drops in percentage points.
The eight insensitive labels were retained from round 1; LTS, T, CO, and PREC
were reassessed in round 2, which added nobody.

\begin{center}
\small
\begin{tabular}{lrrrrl}
\toprule
 & \multicolumn{2}{c}{Clean coverage} & Wildfire & & \\
Candidate & 90\% interval & 95\% interval & 95\% interval & Drop & Decision \\
\midrule
T     & 82.6 & 89.3 & 87.0 & $+2.3$ & Sensitive \\
CO    & 91.6 & 94.5 & 69.8 & $+24.7$ & Sensitive \\
\midrule
PBLH  & 91.5 & 95.0 & 95.6 & $-0.7$ & Insensitive \\
CFLOW & 88.9 & 93.3 & 96.2 & $-2.9$ & Insensitive \\
CFMID & 91.3 & 96.0 & 95.9 & $+0.1$ & Insensitive \\
SWGDN & 87.6 & 93.9 & 93.5 & $+0.4$ & Insensitive \\
LWP   & 85.7 & 93.1 & 93.6 & $-0.5$ & Insensitive \\
OMEGA & 88.2 & 93.8 & 94.7 & $-0.9$ & Insensitive \\
LWC   & 88.7 & 93.3 & 94.3 & $-1.1$ & Insensitive \\
RH    & 87.4 & 91.4 & 93.1 & $-1.7$ & Insensitive \\
\midrule
LTS   & 68.1 & 83.4 & --- & --- & Undetermined \\
PREC  & 69.2 & 75.9 & --- & --- & Undetermined \\
\bottomrule
\end{tabular}
\end{center}

\smallskip
Drops were calculated before rounding, so subtracting the displayed coverages
can differ by 0.1 point. Names omit the stored feature transforms: PBLH, LWP,
LWC, CO, PREC, DUST, and SO2EM use fifth roots; CHL and DMS use square roots.

\textbf{How much should we trust these labels?}
CO has a 24.7-point drop and clean 95\% coverage of 94.5\%. T's drop is only 2.3
points; its 89.3\% clean coverage passes a wide 88.5--101.5\% band. LTS and PREC
fail calibration, so wildfire coverage is not assessed. The two-point cutoff is
not a significance test, and these labels do not establish causality. Accepting
a variable at arrival also does not establish that its full history is insensitive.

\textbf{Shuffle order results.}
Each shuffle seed permutes candidates once and retains that order among survivors.
All runs use arrival-only targets and training/split seed 2026; split indices match
across completed runs. Only candidate order is deliberately changed.
Status below is as of 23 September 2026.

\begin{center}
\small
\begin{tabular}{llrll}
\toprule
Order & Job & Rounds & Sensitive & Undetermined \\
\midrule
Original & 3110691 & 2 & T, CO & LTS, PREC \\
Seed 2027 & 3110763 & 4 & CO & LTS \\
Seed 2028 & 3112708 & 3 & CO & LTS \\
Seed 2029 & 3112709 & 3 & CO & LTS \\
Seed 2030 & 3147106 & \multicolumn{3}{l}{Pending; no results yet} \\
Seed 2031 & 3147107 & \multicolumn{3}{l}{Pending; no results yet} \\
\bottomrule
\end{tabular}
\end{center}

The three completed shuffles agree: CO sensitive, LTS undetermined, and the other
ten tested variables insensitive (20 including the starting ten). Compared with
the original order, only T and PREC change, both becoming insensitive. Thus order
changes the labels, but does not resolve calibration concerns: T's 83.0\% clean
95\% coverage in seed 2027 passes a wide 80.2--109.8\% band. Pending runs are not
included in the agreement count.

{\footnotesize
\textbf{Files.} Source logs are under \nolinkurl{slurm/wildfire_sensitivity/};
filenames end in \texttt{\_<job ID>.log}, using the IDs above.\\
Models, mappings, and splits: each job's \texttt{arrival/} folder under\\
\nolinkurl{/storage3/fs1/myu/Active/felixhu/weather_checkpoints/wildfire_sensitivity/}.
}

\begin{thebibliography}{1}
\small
\bibitem{kuleshov2018}
V.~Kuleshov, N.~Fenner, and S.~Ermon (2018).
\emph{Accurate Uncertainties for Deep Learning Using Calibrated Regression}.
ICML, PMLR 80, pp.~2796--2804.
\url{https://proceedings.mlr.press/v80/kuleshov18a.html}.
\end{thebibliography}
\end{document}
